
%----------------------------------------------------------------------------------------
%	Baseline Design
%----------------------------------------------------------------------------------------
\section{Baseline Design}

\subsection{Subsystem Selections}
Table \ref{tbl:baseline} shows a summary of the subsystem selections resulting from the trade studies.  Feasibility analysis has been started or outlined for each, and will continue with CDR to get a final selection of a specific product or component.
\begin{table}[H]
\centering
\begin{tabular}{|l|c|c|} 
\hline
    \textbf{Critical Project Element} & \textbf{Component} & \textbf{Selection}
    \\
\hline
    \textbf{Launch} & Material & Aluminum 7075
    \\
\hline
    \textbf{} & Chassis & EnduroSat
    \\
\hline
    \textbf{Control} & GNC & Torque Rods
    \\
\hline
    \textbf{Power} & Solar Panels & Stationary Solar Panel
    \\
\hline 
    \textbf{} & Battery & Lithium Ion
    \\
\hline
    \textbf{Thermal} & Passive & Tapes/Paints
    \\
\hline
    \textbf{} & & Deployable Radiators
    \\
\hline
    \textbf{Communication} & Transmitter & UHF/VHF
    \\
\hline
    \textbf{} & Antenna & Vertical Monopole Array
    \\
\hline
    \textbf{End of Life} &  & Electromagnetic Tether
    \\    
\hline
\end{tabular}
\caption{Baseline design selections}
\label{tbl:baseline}
\end{table}

\subsection{Non-CPE Subsystem Selections}

The following subsystems were not included with the in depth trade study analysis performed on CPEs.  Results from the CDD are highlighted here.  Feasibility analysis will still be required to ensure the entire system functions together.

\subsubsection{Payload} High energy particle detector
%\noindent\textbf{Payload:} High energy particle detector
\begin{itemize}
    \item Mass: 1.25 kg
    \item Volume: \textasciitilde100 cm\textsuperscript{3}, cylindrical envelope \textasciitilde3.2 cm radius, 6 cm length
    \item Thermal Limits: $-20\degree$C to $+30\degree$C
    \item FOV: $~50\degree$
    \item Power: <1 W
\end{itemize}

\subsubsection{Avionics}
COTS Highly Integrated On-board Computing System
\begin{itemize}[topsep=0pt]
    \item Space heritage: Many TRL 9 options
    \item Features: Sensor integration, processor, memory
    \item Attributes: Power consumption, radiation resistance
\end{itemize}

\subsection{CAD Model}
A preliminary CAD model of the STARS CubeSat can be found in Figure \ref{fig:CAD_annotated} below. 
\begin{figure}[H]
    \centering
    \includegraphics[width=0.9\columnwidth]{CAD_annotated_not_4K.jpg}\\
    \caption{Preliminary CAD render of STARS with speculated critical subsystems placement.}
    \label{fig:CAD_annotated} 
\end{figure}  
As the current lack of analysis has precluded specific hardware choices, this is a notional CAD model. Potential zones for critical subsystems have been identified. Their final placement has yet to be determined however, both payload and communications will likely be toward the extremities of the spacecraft, with thermal and power towards the core for optimal heat transfer and efficient wiring.
% I dont think we have the capability to write these subsections
%\subsection{Trade study overview}
%\subsection{Sensitivity Analysis}
%\subsection{Literature Review}
\subsection{Orbital Parameters}
STARS will operate in Low Earth Orbit (LEO) as per \textbf{FR.9} in order to monitor the inner Van Allen radiation belt, which starts at approximately 640 kilometers above the surface of the Earth\cite{VanAllen_1}. As stated in \textbf{DR.9.1} the orbit shall be 500 kilometers sun synchronous, meaning that the satellite is passing over any given location along the orbit's ground track at the same local mean solar time on each pass \cite{Boain}. The specific form of sun synchronous orbit (SSO) selected for this mission is the dawn-dusk orbital pattern. A dawn-dusk orbit is a SSO which has a ground track that follows the line along the Earth that represents the transition between night and day \cite{Boain}. Figure \ref{fig:Dusk_Dawn} provides a 3D render of the proposed dawn-dusk orbit.
\begin{figure}[H]
    \includegraphics[width=0.8\columnwidth]{OrbitSnap.png}\\
    \caption{Dawn-Dusk Sun Synchronous Orbit}
    \label{fig:Dusk_Dawn} 
\end{figure}  
The unique ground track of a dawn-dusk orbit will allow STARS to remain in constant sunlight throughout the entire lifetime of the mission. This will help the power system meet \textbf{DR.5.1} by allowing the solar panels to constantly collect power from the sun. Constant power collection will also allow STARS to stay active for longer periods of time and increase the amount of data collected and transmitted. Remaining in constant sunlight also impacts \textbf{FR.6} which revolves around the temperature and radiation STARS will be experiencing. With constant exposure to sunlight STARS will be under constant thermal loads and solar radiation, which the thermal subsystem will have to account for. This orbit does require not only the specific inclination required to maintain a SSO, but it must also have a specific time of ascending node in order to be considered dawn-dusk. A dawn-dusk orbit must have an ascending node that occurs at either 06:00:00 or 18:00:00. In order for the orbit to 500 kilometers sun synchronous it must have an inclination of approximately 97.4\degree \cite{Garada}. While a dawn dusk orbit presents a few issues for the mission the benefit of constant power generation far outweighs the inconveniences they present.

There is another benefit to this orbital design that originates from the sun synchronous nature of STARS' orbit. Orbital period increases with altitude so the 500 kilometer altitude proposed by \textbf{DR.9.1} will result in over 15 revolutions per day, according to simulations made in STK. Sun synchronous orbits maintain a constant trajectory around the Earth as it rotates. This means that the satellite will make multiple passes over the same area in a single day. Having so many revolutions per day will therefore allow for multiple communication windows in a single day. Multiple communication windows will allow for a greater amount of data transmission and more freedom in the mission design. If STARS needs to go into a state of low power consumption for a few passes then it will still have plenty of time to make the next communication window. Several altitudes were considered for STARS, ranging from 450 to 600 kilometers, with each being analyzed for communication with the ground station over the course of a year. For the simulation the ground station was set to be on top of Engineering Building 3 at North Carolina State University. Table \ref{Altitude} provides the minimum, maximum, mean, and total communication times from STARS to the ground station.  
\begin{table}[H]
\centering
\begin{tabular}{| +c | ^c | ^c | ^c | ^c | }
	\hline
	\rowstyle{\bfseries}%
	Altitude (km) & \shortstack{Minimum Transmit \\ Time (s)} & \shortstack{Maximum Transmit \\ Time (s)} & \shortstack{Mean Transmit \\ Time (s)} & \shortstack{Total Transmit \\ Time (s)} \\ \hline
	450 & 25.743 & 655.946 & 516.093 & 873745 \\ \hline
	500 & 21.610 & 695.943 & 546.680 & 965437 \\ \hline   
	550 & 14.317 & 734.760 & 577.544 & 1056906 \\ \hline
	600 & 24.586 & 772.581 & 607.149 & 1147511 \\ \hline

\end{tabular}
\caption{Communication Windows for Simulated Orbit Altitudes}
\label{Altitude}
\end{table}
This range of orbits was selected because they are the most readily available in ride along programs such as the ISIS Small Satellite Launch Service\cite{ISIS}. Both maximum and mean transmit time increase as altitude increases, which results in higher total transmit times. The mean transmit time for any of these altitudes would fulfill the needs of STARS with regards to data transmission. However higher altitudes are more desirable since higher total transmit times would result in more data transmitted to the ground station. Higher orbits draw closer to the inner Van Allen Belt that STARS is studying, which could present more radiation issues for the CubeSat. The proposed 500 kilometer orbit was considered to be a good median for all these factors, and is one of the most readily available altitudes for ride along launches\cite{ISIS} .